%% ============================================================
%%  FPConnect – Trabajo Fin de Grado
%%  Técnico Superior en Desarrollo de Aplicaciones Multiplataforma
%%  Curso 2025–2026
%% ============================================================
\documentclass[12pt, a4paper]{report}

%% ── Codificación y lengua ────────────────────────────────────
\usepackage[utf8]{inputenc}
\usepackage[T1]{fontenc}
\usepackage[spanish]{babel}

%% ── Tipografía ───────────────────────────────────────────────
\usepackage{lmodern}
\usepackage{amsmath}

%% ── Márgenes ─────────────────────────────────────────────────
\usepackage[top=2.5cm, bottom=2.5cm, left=3cm, right=2.5cm]{geometry}

%% ── Interlineado ─────────────────────────────────────────────
\usepackage{setspace}
\onehalfspacing

%% ── Colores ──────────────────────────────────────────────────
\usepackage{xcolor}
\definecolor{fpgreen}{RGB}{34, 139, 34}
\definecolor{fpgray}{RGB}{80, 80, 80}

%% ── Tablas ───────────────────────────────────────────────────
\usepackage{booktabs}
\usepackage{tabularx}
\usepackage{array}

%% ── Gráficos / figuras ───────────────────────────────────────
\usepackage{graphicx}
\usepackage{float}

%% ── Hipervínculos ────────────────────────────────────────────
\usepackage[
  colorlinks = true,
  linkcolor  = fpgreen,
  urlcolor   = fpgreen,
  citecolor  = fpgreen,
  pdftitle   = {FPConnect – TFG DAM},
  pdfauthor  = {Rafael Marcos and Jose Antonio Roda and Angel Amil Cobacho}
]{hyperref}

%% ── Cabeceras y pies de página ───────────────────────────────
\usepackage{fancyhdr}
\pagestyle{fancy}
\fancyhf{}
\fancyhead[L]{\small\textcolor{fpgray}{FPConnect}}
\fancyhead[R]{\small\textcolor{fpgray}{TFG – DAM 2025–2026}}
\fancyfoot[C]{\thepage}
\setlength{\headheight}{14pt}
\renewcommand{\headrulewidth}{0.4pt}

%% ── Estilo de capítulos ──────────────────────────────────────
\usepackage{titlesec}
\titleformat{\chapter}[block]
  {\normalfont\LARGE\bfseries\color{fpgreen}}
  {\thechapter.}{0.8em}{}[\vspace{0.2em}\titlerule]
\titlespacing*{\chapter}{0pt}{30pt}{20pt}

%% ── Código fuente ────────────────────────────────────────────
\usepackage{listings}
\lstset{
  basicstyle   = \ttfamily\small,
  breaklines   = true,
  frame        = single,
  rulecolor    = \color{fpgray},
  keywordstyle = \color{fpgreen}\bfseries,
  commentstyle = \color{fpgray}\itshape,
  stringstyle  = \color{orange},
}

%% ============================================================
\begin{document}

%% ── Portada ─────────────────────────────────────────────────
\begin{titlepage}
  \centering
  \vspace*{1.5cm}

  {\Large\bfseries TRABAJO FIN DE GRADO\\[0.4em]
  Desarrollo de Aplicaciones Multiplataforma}\\[2cm]

  {\Huge\bfseries\color{fpgreen} \(\boldsymbol{\sim}\) FPConnect \(\boldsymbol{\sim}\)}\\[0.8cm]

  {\large Plataforma web para conectar alumnos de FP,\\
  centros educativos y empresas en Andalucía}\\[2.5cm]

  \begin{tabular}{rl}
    \textbf{Alumnos:}    & Rafael Marcos, Jose Antonio Roda, Angel Amil Cobacho \\[0.4em]
    \textbf{Ciclo:}      & Técnico Superior en Desarrollo de Aplicaciones Multiplataforma (DAM) \\[0.4em]
    \textbf{Curso:}      & 2025 – 2026 \\[0.4em]
    \textbf{Repositorio:}& \url{https://github.com/rafamarcoss/FPCONNECT} \\
  \end{tabular}

  \vfill
  {\small Comunidad Autónoma de Andalucía · 2025}
\end{titlepage}

%% ── Índice ──────────────────────────────────────────────────
\tableofcontents
\clearpage

%% ============================================================
\chapter{Resumen}

\textbf{FPConnect} es un proyecto intermodular orientado al desarrollo de una plataforma digital para
conectar a los tres agentes principales de la Formación Profesional: alumnado, centros educativos y
empresas. La solución se plantea como una aplicación web full-stack con arquitectura cliente-servidor,
roles diferenciados y servicios en tiempo real para mejorar la comunicación y la empleabilidad en el
entorno de FP.

La implementación técnica se basa en un frontend con React y Vite, y un backend con Node.js y
Express, utilizando PostgreSQL como sistema gestor de base de datos y Prisma como ORM. La
autenticación se resuelve con JWT y refresh tokens, y la capa de comunicación en tiempo real se apoya
en Socket.io. A nivel de diseño, el sistema aplica una estructura modular por capas (rutas,
controladores, servicios, middlewares y validadores), facilitando la escalabilidad y el mantenimiento.

El proyecto responde a una necesidad real del sector productivo: reducir la distancia entre el talento
técnico en formación y las empresas que necesitan perfiles tecnológicos. Para ello, incorpora
funcionalidades de perfil profesional, conexión entre usuarios, publicación de contenido y gestión de
información de interés para prácticas y empleo.

Además del desarrollo software, el trabajo integra los resultados de aprendizaje de distintos módulos
del ciclo DAM: modelado de datos, servicios en red, diseño de interfaces, control de versiones,
documentación técnica y criterios de sostenibilidad en infraestructura. La memoria recoge metodología,
resultados, limitaciones y líneas futuras, asegurando trazabilidad entre los objetivos definidos y los
artefactos desarrollados.

\vspace{1em}
\begin{table}[H]
  \centering
  \caption{Ficha técnica del proyecto}
  \begin{tabularx}{\textwidth}{lX}
    \toprule
    \textbf{Dato} & \textbf{Información} \\
    \midrule
    Nombre del proyecto   & FPConnect \\
    Tipo de aplicación    & Web full-stack con backend propio \\
    Ámbito geográfico     & Comunidad Autónoma de Andalucía \\
    Usuarios objetivo     & Alumnos FP, centros educativos, empresas \\
    Stack principal       & React + Vite + Node.js + Express + PostgreSQL + Prisma + Socket.io \\
    Año de desarrollo     & 2025 – 2026 \\
    \bottomrule
  \end{tabularx}
\end{table}

%% ============================================================
\chapter{Glosario (ES/EN)}

\begin{table}[H]
  \centering
  \caption{Términos clave del proyecto}
  \begin{tabularx}{\textwidth}{lXX}
    \toprule
    \textbf{ES} & \textbf{EN} & \textbf{Definición breve} \\
    \midrule
    Autenticación & Authentication & Verificación de identidad del usuario. \\
    Autorización & Authorization & Control de permisos según rol. \\
    ORM & Object-Relational Mapping & Mapeo entre objetos de código y tablas relacionales. \\
    Migración & Migration & Cambio versionado del esquema de base de datos. \\
    Punto final API & API endpoint & Ruta HTTP para consumir una funcionalidad del backend. \\
    Tiempo real & Real-time communication & Intercambio de eventos sin refresco manual. \\
    Despliegue & Deployment & Puesta en producción del sistema. \\
    \bottomrule
  \end{tabularx}
\end{table}

%% ============================================================
\chapter{Planteamiento del problema y justificación}

\section{Contexto del problema}

La Formación Profesional en España, y en Andalucía en particular, ha experimentado un crecimiento
notable en los últimos años tanto en número de alumnos matriculados como en la diversificación de
ciclos formativos ofertados. Sin embargo, la transición del mundo educativo al mercado laboral sigue
siendo, en muchos casos, un proceso poco transparente y con escasos canales formales de comunicación
entre las partes implicadas.

Los alumnos de ciclos como DAM, DAW, ASIR o SMR cuentan con un perfil técnico muy demandado, pero
rara vez disponen de una plataforma especializada donde visibilizar su trabajo y conectar directamente
con empresas de su entorno. LinkedIn, aunque generalista, no está optimizado para el perfil de un
alumno en prácticas o recién graduado de FP. Las bolsas de empleo genéricas tampoco cubren bien
las particularidades del sistema de Formación Profesional.

\section{Motivación del proyecto}

La motivación principal detrás de FPConnect nace de la experiencia propia vivida a lo largo del ciclo
de DAM: la dificultad de encontrar empresas para realizar las prácticas del módulo de FCT, la falta de
visibilidad que tienen los proyectos desarrollados durante el ciclo, y la ausencia de un directorio
actualizado de centros FP en Andalucía con información detallada sobre los ciclos que imparten.

Este proyecto intenta ser una respuesta práctica y funcional a esas carencias, aplicando los
conocimientos adquiridos durante el ciclo en tecnologías actuales y ampliamente demandadas en el
mercado.

%% ============================================================
\chapter{Fundamentación teórica y estado del arte}

El desarrollo de plataformas de empleabilidad educativa se apoya en tres pilares: sistemas de
información en red, persistencia de datos relacionales y diseño centrado en usuario. En los últimos
años, el ecosistema JavaScript ha consolidado arquitecturas desacopladas frontend-backend, donde
React permite construir interfaces componibles y Express facilita la exposición de servicios REST.

En persistencia, PostgreSQL se mantiene como una de las opciones más robustas para aplicaciones
transaccionales, mientras que Prisma simplifica la trazabilidad entre modelo lógico y consultas de
aplicación. En comunicación en tiempo real, Socket.io aporta una capa de eventos bidireccionales
adecuada para mensajería y notificaciones, resolviendo limitaciones de polling tradicional.

Desde la perspectiva funcional, las plataformas generalistas de empleo no cubren completamente las
necesidades del contexto FP (prácticas, convenios centro-empresa, visibilidad de proyectos académicos),
por lo que existe oportunidad para una solución específica con segmentación por rol y ciclo formativo.

%% ============================================================
\chapter{Objetivos}

\begin{enumerate}
  \item Diseñar e implementar una plataforma web funcional con tres roles diferenciados: alumnado,
    centro educativo y empresa.
  \item Modelar una base de datos relacional coherente con los procesos de la aplicación, aplicando
    criterios de integridad y normalización.
  \item Implementar un backend con API REST, autenticación segura y comunicación en tiempo real.
  \item Construir una interfaz usable y \emph{responsive}, con navegación coherente y componentes
    reutilizables.
  \item Documentar técnicamente el proyecto y su cumplimiento intermodular, incluyendo límites,
    riesgos y líneas de mejora.
\end{enumerate}

%% ============================================================
\chapter{Planificación temporal y evaluación de coste}

\section{Planificación temporal}

La ejecución del proyecto se ha organizado en fases iterativas para reducir riesgo técnico y facilitar
la validación continua con el tutor. Cada fase incluye objetivos, entregables y revisión.

\begin{table}[H]
  \centering
  \caption{Cronograma resumido por fases}
  \begin{tabularx}{\textwidth}{lXX}
    \toprule
    \textbf{Fase} & \textbf{Contenido principal} & \textbf{Entregables} \\
    \midrule
    Fase 1 & Análisis inicial, alcance y definición de arquitectura & Documento de alcance y stack seleccionado \\
    \addlinespace
    Fase 2 & Modelado de datos relacional y estructura de backend & Esquema Prisma, migraciones y endpoints base \\
    \addlinespace
    Fase 3 & Implementación frontend por roles y navegación & Vistas por rol y componentes reutilizables \\
    \addlinespace
    Fase 4 & Integración full-stack y comunicación en tiempo real & API conectada, sockets funcionales y pruebas manuales \\
    \addlinespace
    Fase 5 & Ajustes finales, documentación y preparación de defensa & Memoria, presentación y anexos de cumplimiento \\
    \bottomrule
  \end{tabularx}
\end{table}

\section{Evaluación de coste (estimación académica)}

El proyecto se ha desarrollado en contexto académico sin coste de licencias propietarias obligatorias,
priorizando herramientas \emph{open source}. La siguiente estimación representa un escenario de
producción básico para un piloto real.

\begin{table}[H]
  \centering
  \caption{Estimación de coste mensual de operación}
  \begin{tabularx}{\textwidth}{lXr}
    \toprule
    \textbf{Concepto} & \textbf{Descripción} & \textbf{Coste estimado/mes} \\
    \midrule
    VPS backend & Servidor Linux con 2 vCPU, 4 GB RAM & 20 - 35 EUR \\
    \addlinespace
    Base de datos & PostgreSQL gestionado o autoalojado en VPS & 0 - 30 EUR \\
    \addlinespace
    Dominio + TLS & Dominio anual prorrateado + certificados Let's Encrypt & 1 - 2 EUR \\
    \addlinespace
    Backups & Almacenamiento incremental para copias diarias & 3 - 8 EUR \\
    \addlinespace
    Monitorización & Logs y alertas básicas & 0 - 10 EUR \\
    \midrule
    Total estimado & Operación mínima de piloto & 24 - 85 EUR \\
    \bottomrule
  \end{tabularx}
\end{table}

%% ============================================================
\chapter{Análisis de requisitos}

\section{Requisitos funcionales priorizados}

La priorización se ha realizado por impacto en valor de usuario y dependencia técnica. Se distinguen
requisitos de tipo \emph{must-have} (imprescindibles para operación) y \emph{should-have}
(evolutivos de alto valor).

\begin{table}[H]
  \centering
  \caption{Priorización funcional}
  \begin{tabularx}{\textwidth}{lXl}
    \toprule
    \textbf{Código} & \textbf{Requisito} & \textbf{Prioridad} \\
    \midrule
    RF-01 & Registro e inicio de sesión con roles diferenciados & Must-have \\
    RF-02 & Gestión de perfil y datos básicos por tipo de usuario & Must-have \\
    RF-03 & Publicaciones, comentarios e interacciones sociales & Must-have \\
    RF-04 & Conexiones entre usuarios y visibilidad por rol & Should-have \\
    RF-05 & Mensajería y eventos en tiempo real & Should-have \\
    RF-06 & Recuperación personalizada de contraseña & Should-have \\
    \bottomrule
  \end{tabularx}
\end{table}

\section{Requisitos técnicos de seguridad y acceso}

\begin{itemize}
  \item Validación estricta de entrada en backend (Joi) para reducir errores de formato y abuso.
  \item Autenticación con JWT y renovación controlada con refresh token.
  \item Contraseñas protegidas mediante hash con bcrypt y política mínima de complejidad.
  \item Gestión de errores centralizada para no exponer información sensible del sistema.
  \item Control de acceso por rol en rutas privadas y acciones sensibles.
\end{itemize}

%% ============================================================
\chapter{Análisis del Sistema}

\section{Identificación de actores}

Se han identificado tres tipos de usuarios principales que interactúan con la plataforma, cada uno con
necesidades y funcionalidades específicas:

\vspace{0.5em}
\begin{table}[H]
  \centering
  \caption{Actores del sistema}
  \begin{tabularx}{\textwidth}{lXX}
    \toprule
    \textbf{Actor} & \textbf{Descripción} & \textbf{Necesidades principales} \\
    \midrule
    Alumno    & Estudiante de un ciclo formativo de FP en Andalucía
              & Visibilizar su perfil, encontrar empresas y centros, acceder a noticias del sector \\
    \addlinespace
    Centro FP & Instituto o centro educativo que imparte ciclos formativos
              & Gestionar su comunidad de alumnos, conectar con empresas para convenios de FCT \\
    \addlinespace
    Empresa   & Empresa que busca incorporar perfiles técnicos de FP
              & Buscar y filtrar talento, publicar vacantes, establecer contacto con centros \\
    \bottomrule
  \end{tabularx}
\end{table}

\section{Requisitos funcionales}

\subsection{Módulo de autenticación y gestión de roles}

\begin{itemize}
  \item El sistema debe permitir el registro e inicio de sesión de usuarios diferenciando entre los
        tres roles disponibles.
  \item Cada rol debe redirigir al usuario a su \emph{dashboard} correspondiente tras la autenticación.
  \item Las rutas y funcionalidades deben estar protegidas según el rol del usuario autenticado.
\end{itemize}

\subsection{Funcionalidades del Alumno}

\begin{itemize}
  \item Acceso a un \emph{dashboard} personal con estadísticas de perfil y noticias del sector.
  \item Directorio de centros FP con filtros por ciclo formativo y provincia.
  \item Visualización de centros y empresas sobre un mapa interactivo de Andalucía.
  \item Acceso a perfiles públicos de otros alumnos de la plataforma.
  \item Perfil público propio con información de CV, habilidades técnicas, proyectos y experiencia.
  \item \emph{Feed} de noticias sobre becas, convocatorias y eventos relevantes para FP.
\end{itemize}

\subsection{Funcionalidades del Centro FP}

\begin{itemize}
  \item \emph{Dashboard} con métricas del centro: número de alumnos registrados, empresas
        colaboradoras, ciclos impartidos.
  \item Listado y gestión de alumnos registrados en el centro.
  \item Visualización de empresas colaboradoras actuales y potenciales.
  \item Perfil público del centro con información de contacto, ciclos y ubicación.
\end{itemize}

\subsection{Funcionalidades de la Empresa}

\begin{itemize}
  \item Buscador de talento con filtros por ciclo formativo, ciudad y disponibilidad.
  \item Acceso al perfil completo de cada alumno, incluyendo \emph{skills} y proyectos.
  \item Gestión de vacantes de prácticas (FCT) y de empleo.
  \item Directorio de centros FP para establecer convenios de colaboración.
  \item Perfil público de empresa con información corporativa.
\end{itemize}

\section{Requisitos no funcionales}

\begin{table}[H]
  \centering
  \caption{Requisitos no funcionales}
  \begin{tabularx}{\textwidth}{lX}
    \toprule
    \textbf{Requisito} & \textbf{Descripción} \\
    \midrule
    Usabilidad
      & La interfaz debe ser intuitiva y consistente. Un usuario nuevo debe poder orientarse sin
        manual. \\
    \addlinespace
    Rendimiento
      & La aplicación web debe cargar en menos de 3 segundos en condiciones normales de red. \\
    \addlinespace
    Multiplataforma
      & El mismo código base debe funcionar como aplicación web y como aplicación de escritorio
        nativa. \\
    \addlinespace
    Escalabilidad
      & La arquitectura backend basada en Node.js, PostgreSQL y Prisma debe soportar el crecimiento
        en número de usuarios sin cambios estructurales en el frontend. \\
    \addlinespace
    Seguridad
      & La autenticación debe basarse en JWT. Los datos sensibles no deben exponerse en el
        frontend. \\
    \addlinespace
    Mantenibilidad
      & El código debe seguir una estructura modular y estar organizado por componentes
        reutilizables. \\
    \bottomrule
  \end{tabularx}
\end{table}

%% ============================================================
\chapter{Diseño del Sistema}

\section{Arquitectura general}

FPConnect sigue una arquitectura cliente-servidor desacoplada. El frontend, desarrollado con React y
Vite, consume una API REST construida con Node.js y Express. El backend integra autenticación por JWT,
validación de entrada, acceso a datos mediante Prisma ORM sobre PostgreSQL y servicios en tiempo real
con Socket.io.

La estructura del backend se organiza por capas (rutas, controladores, servicios, middlewares,
validadores), lo que permite separar responsabilidades y facilitar el mantenimiento evolutivo.

\vspace{0.5em}
\begin{table}[H]
  \centering
  \caption{Capas tecnológicas de la arquitectura}
  \begin{tabularx}{\textwidth}{llX}
    \toprule
    \textbf{Capa} & \textbf{Tecnología} & \textbf{Responsabilidad} \\
    \midrule
    Frontend web    & React + Vite             & Interfaz de usuario, navegación y estado cliente \\
    Backend API     & Node.js + Express        & Lógica de negocio y exposición de servicios REST \\
    Datos           & PostgreSQL + Prisma ORM  & Persistencia, relaciones y migraciones \\
    Tiempo real     & Socket.io                & Eventos y mensajería bidireccional \\
    Mapas           & Leaflet.js + OpenStreetMap & Visualización geográfica de centros y empresas \\
    Despliegue      & Linux Server + Nginx + PM2 & Ejecución de backend en entorno productivo \\
    \bottomrule
  \end{tabularx}
\end{table}

\section{Estructura de directorios}

El proyecto sigue una organización modular basada en la separación por responsabilidades:

\vspace{0.5em}
\begin{table}[H]
  \centering
  \caption{Estructura de directorios del proyecto}
  \begin{tabularx}{\textwidth}{lX}
    \toprule
    \textbf{Ruta} & \textbf{Contenido} \\
    \midrule
    \texttt{src/components/UI.jsx}   & Componentes reutilizables: Badge, Avatar, Card, NavBar, etc. \\
    \texttt{src/data/datos.js}       & Datos de prueba realistas: centros, alumnos, empresas, noticias \\
    \texttt{src/pages/LoginPage.jsx} & Pantalla de autenticación con selección de rol \\
    \texttt{src/pages/AlumnoApp.jsx} & Dashboard completo del perfil de alumno \\
    \texttt{src/pages/CentroApp.jsx} & Dashboard del centro educativo \\
    \texttt{src/pages/EmpresaApp.jsx}& Dashboard de la empresa con buscador de talento \\
    \texttt{src/styles/global.css}   & Reset CSS y estilos globales compartidos \\
    \texttt{src/App.jsx}             & Router principal con enrutado por rol \\
    \texttt{src/main.jsx}            & Punto de entrada de la aplicación React \\
    \texttt{vite.config.js}          & Configuración del bundler Vite \\
    \bottomrule
  \end{tabularx}
\end{table}

\section{Modelo de datos}

El modelo de datos está diseñado para su implementación en PostgreSQL mediante Prisma ORM. Se han
identificado las siguientes entidades principales y sus relaciones:

\vspace{0.5em}
\begin{table}[H]
  \centering
  \caption{Entidades del modelo de datos}
  \begin{tabularx}{\textwidth}{lXX}
    \toprule
    \textbf{Tabla} & \textbf{Campos principales} & \textbf{Relaciones} \\
    \midrule
    \texttt{usuarios}
      & id, email, rol, created\_at
      & Base de todos los perfiles \\
    \addlinespace
    \texttt{alumnos}
      & id, usuario\_id, nombre, ciclo, centro\_id, ciudad, bio, nota, disponible, cv\_url, foto\_url
      & N:1 con usuarios, N:1 con centros, N:M con skills \\
    \addlinespace
    \texttt{centros}
      & id, usuario\_id, nombre, ciudad, provincia, descripcion, telefono, email, web, lat, lng
      & N:1 con usuarios, N:M con ciclos \\
    \addlinespace
    \texttt{empresas}
      & id, usuario\_id, nombre, sector, ciudad, descripcion, empleados, logo\_url
      & N:1 con usuarios \\
    \addlinespace
    \texttt{vacantes}
      & id, empresa\_id, titulo, tipo, ciclo, estado, created\_at
      & N:1 con empresas \\
    \addlinespace
    \texttt{noticias}
      & id, titulo, categoria, resumen, fecha, urgente, url
      & Entidad independiente \\
    \addlinespace
    \texttt{intereses}
      & id, empresa\_id, alumno\_id, created\_at
      & N:M entre empresas y alumnos \\
    \addlinespace
    \texttt{centros\_ciclos}
      & centro\_id, ciclo
      & N:M entre centros y ciclos \\
    \addlinespace
    \texttt{alumnos\_skills}
      & alumno\_id, skill
      & N:M entre alumnos y habilidades \\
    \bottomrule
  \end{tabularx}
\end{table}

\section{Diseño de la interfaz de usuario}

La interfaz de FPConnect sigue los principios de diseño limpio y moderno, con una paleta de colores
basada en verdes y blancos que transmite frescura y profesionalidad. Se han aplicado los siguientes
criterios de diseño:

\begin{itemize}
  \item \textbf{Consistencia visual:} todos los \emph{dashboards} comparten los mismos componentes
        base (Card, Badge, Avatar, NavBar) adaptados a cada rol.
  \item \textbf{Responsive design:} la aplicación está diseñada para funcionar correctamente tanto en
        escritorio como en dispositivos móviles.
  \item \textbf{Tipografía:} se utiliza \emph{DM Sans} de Google Fonts para toda la interfaz, una
        fuente moderna y de alta legibilidad.
  \item \textbf{Componentes reutilizables:} la arquitectura de componentes en React permite que
        elementos como las tarjetas de perfil o las insignias de ciclo sean reutilizados en todos los
        \emph{dashboards}.
\end{itemize}

%% ============================================================
\chapter{Metodología}

La metodología de trabajo ha combinado desarrollo incremental y validación continua, organizada en
iteraciones cortas:

\begin{enumerate}
  \item Definición de alcance funcional y diseño preliminar de datos.
  \item Implementación de base de proyecto (frontend, backend, autenticación y estructura modular).
  \item Desarrollo de funcionalidades por dominio (usuarios, publicaciones, comentarios y conexiones).
  \item Integración de comunicación en tiempo real y pruebas manuales de flujos críticos.
  \item Documentación técnica y mapeo de cumplimiento respecto a criterios intermodulares.
\end{enumerate}

Para el control de calidad se han aplicado validaciones de entrada, separación de lógica por servicios,
uso de migraciones versionadas y comprobaciones funcionales en entorno local.

%% ============================================================
\chapter{Codificación}

\section{Entorno de programación y herramientas}

El desarrollo se ha realizado en Visual Studio Code con control de versiones Git, Node.js LTS y
gestión de dependencias mediante npm. Para persistencia se utiliza Prisma sobre PostgreSQL y para
comunicación en tiempo real Socket.io.

\section{Aspectos relevantes de implementación}

\subsection{Validación de datos}

Toda petición de entrada en rutas críticas se valida en middleware con esquemas explícitos,
interrumpiendo el flujo cuando la carga útil no cumple formato o reglas de negocio básicas.

\subsection{Control de acceso}

Las rutas privadas exigen token válido y, en funcionalidades concretas, se verifica rol y propiedad del
recurso para evitar accesos no autorizados.

\subsection{Protección de la información}

No se almacenan contraseñas en claro. Los secretos de entorno se externalizan en variables de
configuración y no se incluyen en repositorio público.

\subsection{Estructura modular del backend}

La lógica se reparte por responsabilidades: rutas para entrada HTTP, controladores para coordinación,
servicios para reglas de negocio y acceso a datos, y middlewares para capas transversales.

%% ============================================================
\chapter{Resultados}

\section{Estado actual del desarrollo}

En la versión actual del repositorio se encuentra implementada una base funcional del sistema con
autenticación, backend operativo, persistencia relacional y frontend por roles. Se han validado los
flujos principales de acceso, consulta y publicación.

\vspace{0.5em}
\begin{table}[H]
  \centering
  \caption{Estado de implementación – Fase 1}
  \begin{tabularx}{\textwidth}{lX}
    \toprule
    \textbf{Funcionalidad} & \textbf{Descripción} \\
    \midrule
    Login multi-rol
      & Pantalla de inicio de sesión con selección de rol (alumno, centro, empresa). \\
    \addlinespace
    Dashboard alumno
      & Estadísticas, noticias, directorio de centros con filtros, mapa, comunidad y perfil. \\
    \addlinespace
    Dashboard centro
      & Métricas, gestión de alumnos, empresas colaboradoras y perfil del centro. \\
    \addlinespace
    Dashboard empresa
      & Buscador de talento con filtros, gestión de vacantes y directorio de centros. \\
    \addlinespace
    Datos y persistencia
      & Modelo de datos relacional con Prisma y migraciones versionadas. \\
    \addlinespace
    Backend API
      & API REST con Express, validación Joi y autenticación JWT + refresh token. \\
    \addlinespace
    Mapa interactivo
      & Geolocalización real de centros y empresas con Leaflet.js. \\
    \addlinespace
    Mensajería
      & Sistema de mensajería entre empresas y alumnos. \\
    \addlinespace
    Recuperación de contraseña
      & Flujo personalizado de recuperación con pregunta/respuesta de seguridad. \\
    \addlinespace
    Tiempo real
      & Comunicación simultánea cliente-servidor mediante Socket.io. \\
    \bottomrule
  \end{tabularx}
\end{table}

%% ============================================================
\chapter{Pruebas de ejecución}

\section{Pruebas funcionales}

Se han probado manualmente los recorridos críticos: autenticación, navegación por rol, publicaciones,
comentarios, conexiones y recuperación de contraseña. También se validó respuesta de endpoints con
peticiones HTTP controladas.

\section{Pruebas de usabilidad}

La navegación se revisó para minimizar fricción en tareas frecuentes: iniciar sesión, localizar
información principal, publicar contenido y consultar perfiles.

\section{Pruebas de accesibilidad}

Se ha mantenido contraste visual adecuado en elementos principales, tipografía legible y estructura de
interacción consistente. Como línea de mejora se plantea auditoría formal con herramientas
especializadas (por ejemplo TAW o Lighthouse).

\section{Pruebas de seguridad}

Se verificaron casos básicos de denegación por credenciales inválidas, tokens caducados y payloads no
válidos. La recuperación de contraseña exige verificación previa y actualización segura del hash.

\section{Pruebas de carga (estado actual)}

No se han ejecutado pruebas de carga automatizadas extensivas en esta fase. Se incluye como mejora
prioritaria incorporar escenarios de concurrencia y umbrales de rendimiento medibles.

%% ============================================================
\chapter{Manuales de usuario}

\section{Objetivo}

Facilitar a usuarios no técnicos una guía mínima de operación para comenzar a usar la plataforma y
resolver incidencias comunes.

\section{Requisitos e instalación}

\begin{itemize}
  \item Frontend: \texttt{npm install} y \texttt{npm run dev} en la raíz del proyecto.
  \item Backend: \texttt{npm install} en \texttt{backend/}, base de datos PostgreSQL operativa y
        \texttt{npm run dev}.
  \item Navegador moderno actualizado para uso del cliente web.
\end{itemize}

\section{Operaciones habituales}

\begin{enumerate}
  \item Registro o inicio de sesión con credenciales válidas.
  \item Selección de rol y acceso al panel correspondiente.
  \item Gestión de perfil y consulta de información por dominio (alumno, centro o empresa).
  \item Interacción social mediante publicaciones, comentarios y conexiones.
\end{enumerate}

\section{Mensajes de error comunes}

\begin{itemize}
  \item Credenciales inválidas: revisar email/contraseña o usar recuperación.
  \item Token expirado: repetir autenticación o renovar sesión.
  \item Error de validación: corregir campos marcados y reenviar formulario.
\end{itemize}

%% ============================================================
\chapter{Conclusiones}

\section{Grado de consecución de objetivos}

El proyecto alcanza de forma satisfactoria los objetivos principales definidos: arquitectura full-stack
operativa, modelo relacional implementado, servicios de autenticación y capa de interfaz por roles.
Además, se han integrado funcionalidades de valor añadido como comunicación en tiempo real y
recuperación de contraseña personalizada.

\section{Limitaciones}

Las limitaciones más relevantes han sido el tiempo disponible para ampliar cobertura de pruebas
automatizadas y la falta de empaquetado móvil nativo en esta fase. También existen mejoras pendientes
en observabilidad avanzada y escalado de infraestructura para alta concurrencia.

\section{Futuras líneas de trabajo}

\begin{itemize}
  \item Implementar cliente móvil nativo (React Native o Flutter) reutilizando la API existente.
  \item Incorporar analítica de uso y cuadros de mando para centros y empresas.
  \item Añadir CI/CD completo con ejecución automática de tests y despliegue controlado.
  \item Fortalecer seguridad con rate limiting y políticas avanzadas de gestión de sesiones.
\end{itemize}

%% ============================================================
\chapter{Bibliografía}

\begin{itemize}
  \item PostgreSQL Global Development Group. (2025). \emph{PostgreSQL Documentation}. https://www.postgresql.org/docs/
  \item Prisma Labs. (2026). \emph{Prisma ORM Documentation}. https://www.prisma.io/docs
  \item OpenJS Foundation. (2026). \emph{Node.js Documentation}. https://nodejs.org/docs
  \item Express.js. (2026). \emph{Express Web Framework Documentation}. https://expressjs.com/
  \item React Team. (2026). \emph{React Documentation}. https://react.dev/
  \item Socket.IO Authors. (2026). \emph{Socket.IO Documentation}. https://socket.io/docs/
\end{itemize}

%% ============================================================
\chapter{Anexos}

\begin{itemize}
  \item Anexo A: Matriz de cumplimiento intermodular: \path{Requisitos/MATRIZ_CUMPLIMIENTO_INTERMODULAR.md}
  \item Anexo B: Configuración estructurada del proyecto (YAML): \path{Requisitos/config_proyecto.yaml}
  \item Anexo C: Plan de despliegue Linux: \path{Requisitos/despliegue_produccion_linux.md}
  \item Anexo D: Plan de sostenibilidad: \path{Requisitos/plan_sostenibilidad_cpd.md}
  \item Anexo E: Resumen ejecutivo en inglés y glosario técnico bilingüe.
\end{itemize}

\end{document}
